% r03_r1_cumul: the cumulant bound for the step Gaussian, part II: tails, the normalisation of the cutoff, cumulants of bounded variables, the cutoff in the exponent and in cumulants, the connected Wick expansion and the comparison of two Gaussian laws.

\begin{lemma}[tails]\label{rot:lem:R1tail}
Let $\nu=N(h,G)$ on $(\R^n)^X$ with $\|h\|_\infty\le\beta/2$ and $G\le\gamma_0^{-1}$. Then $\nu(\exists x\in X:\ |z_x|>\beta)\le\epsilon(X,\beta)\defeq2n|X|e^{-\gamma_0\beta^2/(8n)}$; since $\bar B(\beta)\subset\beta K$, also $\nu(1-\chi^X_\beta)\le\epsilon(X,\beta)$.
\end{lemma}
\begin{proof}
$|z_x|>\beta$ forces $|\zeta_x|>\beta/2$ for $\zeta\defeq z-h$, hence $|\zeta^a_x|>\beta/(2\sqrt n)$ for some component $a$; $\zeta^a_x\sim N(0,\sigma^2)$ with $\sigma^2\le\gamma_0^{-1}$, $\Pr(|N(0,\sigma^2)|>t)\le2e^{-t^2/(2\sigma^2)}$, and a union bound over the $n|X|$ components gives the claim.
\end{proof}

\begin{lemma}[the normalisation of the cutoff]\label{rot:lem:R1gci}
Let $\mu$ be a centred Gaussian on $(\R^n)^\Lambda$ all of whose component variances are at most $\gamma_0^{-1}$, and let $2e^{-\gamma_0\beta^2/(2n)}\le\frac12$. Then $\mu(\chi_\beta)\ge\exp\big(-4n|\Lambda|e^{-\gamma_0\beta^2/(2n)}\big)$.
\end{lemma}
\begin{proof}
$\{z_x\in\beta K\}\supset\{|z_x|\le\beta\}\supset\bigcap_{a\le n}\{|z^a_x|\le\beta/\sqrt n\}$. By \v{S}id\'ak's inequality (Z.~\v{S}id\'ak, J.\ Amer.\ Statist.\ Assoc.\ 62 (1967), 626--633; C.~G.~Khatri, Ann.\ Math.\ Statist.\ 38 (1967), 1853--1867): for a centred Gaussian vector $(\xi_1,\dots,\xi_m)$ and $t_1,\dots,t_m>0$, $\Pr(\bigcap_i\{|\xi_i|\le t_i\})\ge\prod_i\Pr(|\xi_i|\le t_i)$; this is the case of one slab against a product of slabs of the Gaussian correlation inequality, and it is applied here to the $n|\Lambda|$ components of the real centred vector $z$. Each factor is at least $1-2e^{-\gamma_0\beta^2/(2n)}$, and $\log(1-u)\ge-2u$ for $u\le\frac12$.
\end{proof}

\begin{lemma}[cumulants of bounded variables]\label{rot:lem:R1cum}
For any probability measure $\nu$ and real variables with $|X_j|\le M_j$, $|\kappa_\nu(X_1,\dots,X_k)|\le k!\,2^{k-1}\prod_jM_j$. If $|X|\le M$, then $\big|\log\int e^{X}d\nu-\sum_{k=1}^{\bar n}\kappa_\nu(X;k)/k!\big|\le2^{\bar n}M^{\bar n+1}\le(2M)^{\bar n+1}$.
\end{lemma}
\begin{proof}
In the M\"obius formula $\kappa_\nu(X_1,\dots,X_k)=\sum_\pi(-1)^{|\pi|-1}(|\pi|-1)!\prod_{B\in\pi}\E_\nu\prod_{j\in B}X_j$, over the set partitions $\pi$ of $\{1,\dots,k\}$, each product of moments is at most $\prod_jM_j$, and $\sum_\pi(|\pi|-1)!$ is at most the number of ordered set partitions, which is at most $k!\,2^{k-1}$ (an ordered partition is a permutation cut into consecutive blocks). For the second claim, $f(s)\defeq\log\int e^{sX}d\nu$ has $f^{(m)}(s)=\kappa_{\nu_s}(X;m)$ with $\nu_s\propto e^{sX}\nu$ a probability measure under which still $|X|\le M$; Taylor's formula at $s=0$ with the Lagrange remainder $f^{(\bar n+1)}(s_*)/(\bar n+1)!$ and the first claim give the bound.
\end{proof}

\begin{lemma}[the cutoff in the exponent]\label{rot:lem:R1ext}
Let $\nu$ be a probability measure, $\chi$ an indicator with $\nu(1-\chi)\le\epsilon$, $W$ measurable with $|W|\le M$, and $e^{2M}\epsilon\le\frac12$. Then \textup{(i)} $\int\chi e^Wd\nu\ge e^{-2e^{2M}\epsilon}\int e^Wd\nu$, and \textup{(ii)} $\int\chi e^Wd\nu\le\int e^{\chi W}d\nu\le e^{2e^M\epsilon}\int\chi e^Wd\nu$.
\end{lemma}
\begin{proof}
(i) $\int(1-\chi)e^W\le e^M\epsilon\le e^{2M}\epsilon\int e^W$, because $\int e^W\ge e^{-M}$; and $1-u\ge e^{-2u}$ for $u\le\frac12$. (ii) On $\{\chi=0\}$, $e^{\chi W}=1\ge0=\chi e^W$, so $\int e^{\chi W}=\int\chi e^W+\nu(1-\chi)$; and $\int\chi e^W\ge e^{-M}(1-\epsilon)\ge e^{-M}/2$, so $\nu(1-\chi)\le2\epsilon e^M\int\chi e^W$.
\end{proof}
The lemma is applied with $W=\Psi\chi$, for which $\chi e^W=\chi e^\Psi$.

\begin{lemma}[connected Wick expansion]\label{rot:lem:R1wick}
Let $\zeta$ be a centred Gaussian vector with covariance $G$ and let $F_j=\zeta^{\alpha_j}$, $j=1,\dots,k$, be monomials, each a multiset of legs. Then $\kappa(F_1,\dots,F_k)=\sum_\pi\prod_{\{u,v\}\in\pi}G(u,v)$, the sum over the pairings $\pi$ of all legs that connect the $k$ factors (the graph on $\{1,\dots,k\}$ with an edge for each pair joining legs of two different factors is connected; pairs inside a factor are allowed). For $N(h,G)$ the same holds after the shift $z=h+\zeta$, the constants dropping out of joint cumulants with $k\ge2$.
\end{lemma}
\begin{proof}
Isserlis' formula gives $\E\prod_{j\in B}F_j=\sum\prod G$ over all pairings of the legs of the factors in $B$. Grouping each pairing by the partition of $B$ into the connected components of its factor graph gives $\E\prod_{j\in B}F_j=\sum_\sigma\prod_{b\in\sigma}\mathsf C(b)$, with $\mathsf C(b)$ the sum over the connected pairings of the legs of $b$; the moment--cumulant relation $\E\prod_{j\in B}F_j=\sum_\sigma\prod_{b\in\sigma}\kappa(b)$ determines the $\kappa(b)$ from the moments recursively, so $\kappa=\mathsf C$.
\end{proof}

\begin{lemma}[anchored connected bound]\label{rot:lem:R1conn}
Let $\nu=N(h,G)$ on $(\R^n)^X$ with $\|G(x,y)\|\le g_0e^{-\theta|x-y|}$ and $\|h\|_\infty\le\beta$. Let $F_1,\dots,F_k$ ($k\ge2$) be polynomials of degree at most $D$ in the variables of $X$, sorted by supports as in (R3), with $\|F_j\|_{\theta'}\le a_j$ for some $\theta'\le\theta$, and let $S_1,S_2\subset\mathbb{L}$ be such that every monomial of $F_1$ has support meeting $S_1$ and every monomial of $F_2$ has support inside $S_2$. Then
\[
|\kappa_\nu(F_1,\dots,F_k)|\le|S_1\cap X|\,K_k(\theta')\,(1+\beta)^{kD}\prod_ja_j\,e^{-(\theta'/2)\dist(S_1,S_2)},\qquad K_k(\theta')\defeq(kD)^{kD/2}\max\{1,g_0\}^{kD/2}\,\tilde c(\theta'/2)^{k-1},
\]
and the same bound without the distance factor when no $S_2$ is given (the anchored bound). For $k=1$, $|\E_\nu F_1|\le|S_1\cap X|\,D^{D/2}\max\{1,g_0\}^{D/2}(1+\beta)^Da_1$. For centred $\nu$ the factor $(1+\beta)^{kD}$ is absent.
\end{lemma}
\begin{proof}
(1) Shift $z=h+\zeta$: a monomial $z^\alpha$ with support $Y$ becomes $\sum_{\alpha'\le\alpha}\binom{\alpha}{\alpha'}h^{\alpha-\alpha'}\zeta^{\alpha'}$, whose coefficients have modulus sum at most $\prod(1+|h|)^{\alpha}\le(1+\beta)^D$; the pieces are kept grouped under the original support $Y$ (their sites lie in $Y$), so the shifted $F'_j$ has $\|F'_j\|_{\theta'}\le(1+\beta)^Da_j$, and the constant pieces drop out for $k\ge2$. For $h=0$ there is no shift and no factor. (2) By \Cref{rot:lem:R1wick}, $|\kappa_\nu(F_1,\dots,F_k)|\le\sum_{Y_1,\dots,Y_k}\sum_{\text{monomials}}\prod|\text{coefficient}|\sum_{\text{connected }\pi}\prod\|G\|$; the factors have at most $kD$ legs in all, so for fixed monomials there are at most $(kD)^{kD/2}$ pairings. (3) For a connected $\pi$ fix a spanning tree of its factor graph, rooted at factor $1$, with $k-1$ tree pairs. Sum the root over the sites $a\in S_1\cap X$ and the supports $Y_1\ni a$: at most $|S_1\cap X|\,\|F'_1\|_{\theta'}$ times the rest. Each child factor $j$ is entered through a tree pair $(y,y')$ with $y$ in the parent's support and $y'\in Y_j$, which gives $\sum_{y'}g_0e^{-\theta|y-y'|}\sum_{Y_j\ni y'}|F'_{j,Y_j}|_1$; the pairs not in the tree are bounded by $g_0$ each. (4) Write $|F'_{j,Y}|_1=\big(|F'_{j,Y}|_1e^{\theta'\ell(Y)}\big)e^{-\theta'\ell(Y)}$ and $e^{-\theta|y-y'|}\le e^{-(\theta'/2)|y-y'|}e^{-(\theta'/2)|y-y'|}$. The tree path from factor $1$ to factor $2$ lifts to a path in $\mathbb{L}$ from $a\in S_1$ to a point of $Y_2\subset S_2$, made of tree pairs and of paths inside shortest connecting graphs of the traversed supports $Y_j$, each of length at most $\ell(Y_j)$ (\Cref{rot:lem:R1lat}(iii)); hence $e^{-(\theta'/2)(\sum_j\ell(Y_j)+\sum_{\text{tree}}|y-y'|)}\le e^{-(\theta'/2)\dist(S_1,S_2)}$. The remaining halves give $\tilde c(\theta'/2)$ per tree pair (\Cref{rot:lem:R1lat}(ii)) and $\|F'_j\|_{\theta'}$ per factor. For $k=1$ the same count applies without a tree.
\end{proof}

\begin{lemma}[cumulants with and without the cutoff]\label{rot:lem:R1chi}
Let $\nu=N(h,G)$ on $(\R^n)^X$ with $\|h\|_\infty\le\beta/2$, $G\le\gamma_0^{-1}$ and $\beta\ge2$; put $\chi\defeq\chi^X_\beta$ and $\epsilon\defeq\nu(1-\chi)$. Let $F_1,\dots,F_k$ ($k\le\bar n$) be polynomials in the variables of $X$ of degree at most $D$ with $\|F_j\|_0\le a_j$, and assume the moment condition
\begin{equation}\label{rot:eq:R1MC}
2^q\,|X|\,(2n)^{1/2}\big(n^q(2q-1)!!\big)^{1/2}\gamma_0^{-q/2}\,e^{-\gamma_0\beta^2/(4n)}\le(2\beta)^q\qquad(1\le q\le2\bar nD).
\end{equation}
Then $|\kappa_\nu(\chi F_1,\dots,\chi F_k)-\kappa_\nu(F_1,\dots,F_k)|\le k\cdot k!\cdot4^k\prod_j(a_j|X|)\,(2\beta)^{kD}\,\epsilon^{1/2}$.
\end{lemma}
\begin{proof}
Let $Z\defeq\max\{1,\max_{x\in X}|z_x|\}$ and $\zeta\defeq z-h$. If $\max_x|\zeta_x|\le\beta$ then $Z\le2\beta$, otherwise $Z\le2\max_x|\zeta_x|$; so $\E Z^q\le(2\beta)^q+2^q\sum_x\E\big[|\zeta_x|^q\mathbf 1\{|\zeta_x|>\beta\}\big]\le2(2\beta)^q$ by the Cauchy--Schwarz inequality, $\E|\zeta_x|^{2q}\le n^q(2q-1)!!\gamma_0^{-q}$ (H\"older's inequality over the $n$ components and the Gaussian moments), $\Pr(|\zeta_x|>\beta)\le2ne^{-\gamma_0\beta^2/(2n)}$ and \eqref{rot:eq:R1MC}. Since $|F_j|\le a_j|X|Z^D$, every product of moments obeys $|\E\prod_{j\in B}F_j|\le2\prod_{j\in B}(a_j|X|)(2\beta)^{|B|D}$, and $|\E[(1-\chi)\prod_{j\in B}F_j]|\le\epsilon^{1/2}(\E\prod_{j\in B}F_j^2)^{1/2}\le\sqrt2\,\epsilon^{1/2}\prod_{j\in B}(a_j|X|)(2\beta)^{|B|D}$. In the M\"obius formula ($\chi^{|B|}=\chi$ on each block) replace the blocks one at a time: each partition contributes at most $k\cdot\sqrt2\cdot2^{k-1}\epsilon^{1/2}\prod_j(a_j|X|)(2\beta)^{kD}$, and $\sum_\pi(|\pi|-1)!\le k!\,2^{k-1}$.
\end{proof}
The moment condition \eqref{rot:eq:R1MC} is monotone: once it holds for a pair $(|X|,\beta)$ it holds for smaller $|X|$ and larger $\beta$.

\begin{lemma}[comparing two Gaussian laws]\label{rot:lem:R1repl}
Let $\nu,\tilde\nu$ be Gaussian laws of $(z_x)_{x\in S}$, $S$ finite, with means $h,\tilde h$ of sup norm at most $\beta/2$ and covariances $G,\tilde G$ all of whose blocks have norm at most $\gamma_0^{-1}$, and let $\sup_x|h(x)-\tilde h(x)|\le\eta_1$, $\sup_{x,y}\|(G-\tilde G)(x,y)\|\le\eta_2$. For polynomials $F_1,\dots,F_k$ in the variables of $S$ of degree at most $D$ with $\|F_j\|_0\le a_j$,
\[
|\kappa_\nu(F_1,\dots,F_k)-\kappa_{\tilde\nu}(F_1,\dots,F_k)|\le k\cdot k!\cdot4^k\,(kD)^{kD+1}\,(1+\beta+\gamma_0^{-1})^{kD}\prod_j(a_j|S|)\,\max\{\eta_1,\eta_2\}.
\]
\end{lemma}
\begin{proof}
A moment of $L\le kD$ coordinate variables is, by Isserlis' formula for a Gaussian with mean, a sum over at most $L^L$ partial pairings of products of at most $L$ factors (covariance entries or means), each of modulus at most $1+\beta+\gamma_0^{-1}$ under both laws; replacing the factors one at a time, two corresponding products differ by at most $L\max\{\eta_1,\eta_2\}(1+\beta+\gamma_0^{-1})^L$. The coefficient sums contribute $\prod_j(a_j|S|)$, and the M\"obius formula with the telescoping of \Cref{rot:lem:R1chi} gives the bound.
\end{proof}
